
\begin{figure}[!htbp]
  \centering
  \includegraphics[width=\linewidth]{results/figures/step4/step4_comparison.pdf}
  \caption{Qualitative comparison on Cityscapes val: Input, Ground Truth, 
  EoMT-COCO (panoptic prediction, COCO label space), and EoMT-Cityscapes 
  (semantic prediction). EoMT-COCO produces segmentations in a different class space.}
  \label{fig:step4_comparison}
\end{figure}
\section{Results}
\begin{table}[!htbp]
\centering
\small
\caption{Selected per-class IoU (\%) on \texttt{Cityscapes val}. EoMT-COCO scores 0\% on classes with no COCO equivalent.}
\label{tab:perclass}
\begin{tabular}{lcc}
\toprule
Class & EoMT-COCO & EoMT-Cityscapes \\
\midrule
pole          & \textbf{0.00}  & 71.04 \\
rider         & \textbf{0.00}  & 71.11 \\
traffic sign  & \textbf{4.74}  & 82.13 \\
\midrule
road          & 94.13 & 98.40 \\
car           & 84.33 & 95.54 \\
bicycle       & 73.34 & 81.27 \\
\midrule
\textbf{Mean} & \textbf{55.17} & \textbf{81.68} \\
\bottomrule
\end{tabular}
\end{table}
\subsection{Semantic Segmentation}
EoMT-Cityscapes is the best model with 81.68\% mIoU among the compared model. Beside them ERFNet (72.50\%) is competitive for a pixel-based model.
EoMT-COCO scores 0\% IoU on \texttt{pole, rider}, and near-0\% on \texttt{traffic sign} (Table~\ref{tab:perclass}) because COCO has no equivalent label.
It is not the problem of the model but a label space mismatch. It is the motivation for fine-tuning.
Fine-tuning lifts mIoU from 55.17 to 78.85, recovers the most of the gap. This result shows that COCO's DINOv2-backed visual priors work well, the only problem was the classification head, not the features.
The difference is visible in Figure~\ref{fig:step4_comparison} for the class-label mismatch problem.
\subsection{Anomaly Segmentation}

\paragraph{Pixel-based vs mask-based.}

We see these results from the comparison of ERFNet and EoMT (Table~\ref{tab:temperature_combined}):
\begin{itemize}
    \item On RoadAnomaly-21 dataset, ERFNet's MaxLogit is 38.31, and EoMT-Cityscapes' MaxEntropy is 77.88. There is a 40\% point difference
    \item On RoadObstacle-21 dataset, ERFNet's MaxLogit is 4.63 and EoMT-Cityscapes' MSP is 90.11. There is a 85\% point difference
    \item On RoadAnomaly dataset, ERFNet's MaxLogit 15.58 and EoMT-Cityscapes' MaxEntropy 78.14. There is a 63\% point difference
\end{itemize}
With these results, we find that on most benchmarks, the EoMT's models substantially outperform the ERFNet. Mask-based classification 
creates smooth shapes because it looks at the whole object at once. However, pixel-softmax looks at each pixel one by one, that's why it is vulnerable to noise.

\paragraph{Effect of fine-tuning.}

To see the effect of the fine-tuning, we compare the EoMT-Cityscapes with EoMT-Fine-tuned on AuPRC (Table~\ref{tab:temperature_combined}).
\begin{itemize}
    \item On RoadAnomaly-21 dataset: MSP score drops from 74.36\% to 57.15\% (~17\% change)
    \item On RoadObstacle-21 dataset: MSP score drops from 90.11\% to 73.61\% (~16\% change)
    \item On RoadAnomaly dataset: MSP score drops from 75.69\% to 54.99\% (~21\% change)
    \item On FS L\&F dataset: MSP score increases from 25.42\% to 27.11\% (~2\% change)
\end{itemize}
According to these results, EoMT-Fine-tuned recovers closed set mIoU (Table~\ref{tab:temperature_combined}) but decrease the anomaly scores compared to training from scratch on Cityscapes.
This happens because training the model for a short time makes it hyper-focused on normal objects, and this ruins its ability to spot unusual things. 

\paragraph{Comparison of scoring functions.}
Interesting finding here is that MaxLogit beats MSP and MaxEntropy on ERFNet (38.31 vs 29.09 vs 30.96 on RA-21), but on EoMT they are almost same.
This result shows that scoring-method choice depends on the architecture. Figure~\ref{fig:step8} shows score maps for the four methods on RA-21 images.
\textcolor{blue}{After revising RbA, we report it as a rejected-by-all known-class-mass score instead of a temperature-scaled softmax score. For EoMT-COCO, the revised RbA improves the anomaly score on RoadObstacle21 compared with the other EoMT-COCO scoring functions, reaching 36.75\% AuPRC, while MSP, MaxLogit, and MaxEntropy remain below 4\%. On RoadAnomaly21, revised RbA reaches 37.44\% AuPRC, which is competitive with MSP and MaxLogit but still below MaxEntropy. However, RbA still gives high FPR95 on some datasets, especially RoadObstacle21 and FS Static, meaning that it can detect some anomalous regions but may also produce many false positives.}
\begin{figure}[!htbp]
\centering
\includegraphics[width=\linewidth]{results/figures/step8/step8_anomaly_maps.png}
\caption{Qualitative comparison of anomaly score maps on RoadAnomaly21. 
Each row shows an input image alongside score maps produced by MSP, MaxLogit, 
MaxEntropy, and RbA for EoMT-Cityscapes.}
\label{fig:step8}
\end{figure}
%\paragraph{Qualitative results.}
%Figure~\ref{fig:step8} shows anomaly score maps for the four scoring methods on representative RoadAnomaly21 images.

\input{results/tables/latex_tables/combined_temperature_table.tex}

\subsection{Temperature Scaling Study}

For temperature scaling, we tried $T \in \{0.5, 0.75, 1.0, 1.1, 1.5, 2.0\}$ on MSP and MaxEntropy 
for EoMT-Cityscapes, EoMT-Fine-tuned, and ERFNet.
\textcolor{blue}{In the revised version, RbA is excluded from temperature scaling because the final RbA score is computed as the negative total known-class mass over the EoMT semantic score map, not as a softmax-calibrated probability score.}
Main finding that AuPRC is mostly stable across $T$ (Table~\ref{tab:temperature_combined}).
\begin{itemize}
    \item In EoMT-Cityscapes (Table~\ref{tab:temperature_combined}), MSP AuPRC on RA-21 stays in a 0.3\% window, and MaxEntropy stays in a 0.5\% window.
    \item In EoMT-Fine-tuned (Table~\ref{tab:temperature_combined}), It follows similar pattern as EoMT-Cityscapes
    \item In ERFNet (Table~\ref{tab:temperature_combined}), MSP on RA-21 has a change range for 3\%, but its AuPRC is so low that the gain is not meaningful. 
\end{itemize}

\textcolor{blue}{The previous RbA temperature analysis was removed because it no longer corresponds to the revised RbA implementation. In the final version, RbA is reported without temperature scaling.}
As result, we found that T=1.0 is competitive for AuPRC most settings. Temperature scaling doesn't meaningfully improve AuPRC for MSP and MaxEntropy, and it is not a consistent gain.
Post-hoc temperature calibration is not a reliable method for improving anomaly segmentation on these benchmarks.



Interesting finding that RbA's AuPRC is stable, but its FPR95 swings with T:
\begin{itemize}
    \item EoMT-Cityscapes RbA on FS L\&F: FPR95=67.27 at T=0.5 vs 8.79 at T=1.0
    \item EoMT-Cityscapes RbA on RoadAnomaly: FPR95's score decreases from 20.41 at T=1.0 to 17.96 at T=2.0.
    \item EoMT-Fine-tuned RbA on RoadAnomaly: FPR95 = 81.57 at T=0.5 vs 12.09 at T=2.0
\end{itemize}

As result, we found that T=1.0 is competitive for AuPRC most settings. Temperature scaling doesn't meaningfully improve AuPRC, and can shift the FPR95 results, especially for RbA, but it is not a consistent gain.
Post-hoc temperature calibration is not a reliable method for improving anomaly segmentation on these benchmarks.
